% Draft, sections 1-5. Written from results/RESULTS.md and RESEARCH_BRIEF_v3.md.
% Every number is a macro from generated/numbers.tex (src/scripts/paper_results.py);
% this file types no measured value of its own. Compiled on the LaTeX host, never
% on the Jetson. The Interspeech class replaces `article' there.
\documentclass[a4paper,10pt,twocolumn]{article}
\usepackage[utf8]{inputenc}
\usepackage{amsmath,amssymb}
\usepackage{graphicx}
\usepackage{booktabs}
\usepackage[hidelinks]{hyperref}
% Generated by src/scripts/paper_results.py from the claim registry.
% Do not edit.
\newcommand{\PthreeRed}{+357 ms [+294, +412]}  % p3
\newcommand{\PthreeN}{70}  % p3
\newcommand{\DevTTFA}{+2 ms [-15, +15]}  % p3_dev
\newcommand{\SttShareMulti}{62}  % stages_multi_step
\newcommand{\LlmShareMulti}{5}  % stages_multi_step
\newcommand{\TTFAMulti}{4353}  % stages_multi_step
\newcommand{\SttShareDev}{52}  % stages_dev
\newcommand{\SttShareShort}{47}  % stages_short_digit
\newcommand{\SttShareLong}{46}  % stages_long_digit
\newcommand{\ThinkerOverlap}{+96.5 ms [+79.4, +109.8]}  % thinker_overlap
\newcommand{\ThinkerOverlapN}{258}  % thinker_overlap
\newcommand{\ListenerC}{492}  % listener_c
\newcommand{\ListenerCCI}{[439, 574]}  % listener_c
\newcommand{\PartialsDev}{+305.7 ms [+188.0, +452.6]}  % listener_intervention
\newcommand{\TaxBallastA}{+113 ms [+78, +140]}  % tax_ballast_7dc19e
\newcommand{\TaxBallastB}{+249 ms [+169, +275]}  % tax_ballast_ace562
\newcommand{\TaxFullA}{+288 ms [+200, +348]}  % tax_full_37443c
\newcommand{\WordsCommitted}{0.409 [0.340, 0.454]}  % p8_listener
\newcommand{\EnergyListener}{+17.83 J [+17.12, +18.58]}  % energy_listener
\newcommand{\SpecEnergy}{+35.258 [+34.561, +35.659]}  % spec_energy
\newcommand{\SpecTTFA}{+32.375 [-91.583, +92.450]}  % spec_ttfa
\newcommand{\SpecTokens}{96}  % spec_survival
\newcommand{\Fstar}{0.861 [0.853, 0.865]}  % fstar
\newcommand{\PnineTTFA}{+462 ms [+57, +737]}  % p9c
\newcommand{\ResidencyTTFA}{+1.8 ms [-30.3, +37.0]}  % residency
\newcommand{\WorkerPSS}{121.4}  % worker_memory
\newcommand{\PtwoTTFA}{-155.121}  % pass2
\newcommand{\PlannerDP}{+12 ms [+0, +30]}  % planner_dp
\newcommand{\PlannerNoDuty}{+46 ms [+25, +71]}  % planner_noduty
\newcommand{\StaleMean}{1.74}  % stale_net
\newcommand{\StaleN}{58}  % stale
\newcommand{\PfourMean}{+0.0043}  % p4
\newcommand{\PsixRatio}{0.95x [0.94, 0.96]}  % p6
\newcommand{\SevenJms}{+0.0499 J/ms [+0.0425, +0.0607]}  % p7_jms
\newcommand{\ResidencyMem}{19.1}  % residency_memory
\newcommand{\PfourCI}{[-0.0039, +0.0135]}  % p4_ci
\newcommand{\PfourWorse}{14}  % p4_ci
\newcommand{\PfourBetter}{8}  % p4_ci
\newcommand{\FstarUsable}{0.155}  % fstar
\newcommand{\FstarN}{36}  % fstar
\newcommand{\WordsN}{70}  % p8_listener
\newcommand{\StaleTotal}{101.1}  % stale_net
\newcommand{\AloneSTT}{1723 ms [1520, 1866]}  % fig2_alone
\newcommand{\AloneFallback}{0}  % fig2_alone
\newcommand{\AloneHuge}{196}  % fig2_alone
\newcommand{\BallastFallback}{285}  % fig2_ballast
\newcommand{\BallastHuge}{46}  % fig2_ballast
\newcommand{\TaxFullB}{+212 ms [+153, +248]}  % tax_full_74e10a
\newcommand{\PnineRatio}{0.93x [0.92, 0.98]}  % p9c
\newcommand{\PnineN}{16}  % p9c
\newcommand{\PartialsDevN}{15}  % listener_intervention
\newcommand{\ThinkerFinished}{+17.5 ms [-24.4, +76.3]}  % thinker_finished
\newcommand{\ThinkerFinishedN}{66}  % thinker_finished
\newcommand{\SpecN}{22}  % spec_ttfa
\newcommand{\SpecItems}{38}  % spec_survival
\newcommand{\SpecAccepted}{0.000}  % spec_survival
\newcommand{\GreedySeventyFiveK}{5/80}  % probe_greedy_075
\newcommand{\TriggerPoints}{198}  % trigger_points
\newcommand{\TriggerUtts}{16}  % trigger_points
\newcommand{\LlmRate}{31.8}  % llm_rate
\newcommand{\PtwoaPop}{960 turns, 43 runs}  % p2a_pop
\newcommand{\AttrN}{32}  % attr_paired
\newcommand{\ThpMZero}{20}  % thp_m0
\newcommand{\ThpMOne}{44}  % thp_m1
\newcommand{\ThpMTwo}{46}  % thp_m2
\newcommand{\ThpMThree}{66}  % thp_m3
\newcommand{\PgFirst}{0/15}  % pg_first
\newcommand{\PgSurvival}{2.4\%}  % pg_survival
\newcommand{\PgSurvivalCI}{[0.0\%, 7.5\%]}  % pg_survival
\newcommand{\PgSurvivalN}{16}  % pg_survival
\newcommand{\PgBlockedN}{48}  % pg_blocked
\newcommand{\CarryN}{12}  % carry
\newcommand{\PtwoSixN}{211}  % p6_pass2
\newcommand{\PtwoDevN}{16}  % pass2_dev
\newcommand{\OnsetNinetySixRun}{+205.1 [+88.9, +295.5]}  % onset_c96
\newcommand{\OnsetNinetySixRunN}{16}  % onset_c96
\newcommand{\OnsetSixtyFourFin}{+128.8 [+57.6, +155.7]}  % onset_c64
\newcommand{\OnsetSixtyFourFinN}{29}  % onset_c64
\newcommand{\LeadOneN}{12}  % lead_one
\newcommand{\LeadSharedN}{30}  % lead_shared
\newcommand{\LeadNewN}{18}  % lead_new
\newcommand{\EpaWall}{11.53}  % epa_wall
\newcommand{\EpaN}{200}  % epa_wall
\newcommand{\EpaRtf}{4.16}  % epa_rtf
\newcommand{\ResidencyFull}{+12.1 ms [-31.1, +39.1]}  % residency_full
\newcommand{\TriggerUnindexed}{6}  % trigger_points
\newcommand{\PlannerVal}{+399 ms [+362, +477]}  % planner_validation
\newcommand{\PtwoN}{62}  % pass2
\newcommand{\ResidencyMemRange}{+14.0 to +43.1}  % residency_memory
\newcommand{\AttrPaired}{+2.6 ms [-99.5, +182.3]}  % attr_paired
\newcommand{\PgBlocked}{+114 ms [+61, +161]}  % pg_blocked
\newcommand{\Carry}{+169 ms [+65, +1204]}  % carry
\newcommand{\PtwoSix}{1.31x [1.16, 1.46]}  % p6_pass2
\newcommand{\PtwoDev}{-395 ms [-527, -253]}  % pass2_dev
\newcommand{\PnineKill}{2.06x [2.01, 2.07]}  % p9_kill
\newcommand{\LeadOne}{897 ms [607, 1096]}  % lead_one
\newcommand{\LeadShared}{561 ms [502, 1351]}  % lead_shared
\newcommand{\LeadSharedCS}{1685 ms [1369, 1908]}  % lead_cs_shared
\newcommand{\LeadNew}{477 ms [276, 543]}  % lead_new
\newcommand{\GreedySeventyFive}{0.062 [0.027, 0.138]}  % probe_greedy_075
\newcommand{\GreedyNinety}{0.188 [0.117, 0.287]}  % probe_greedy_090
\newcommand{\GreedyFifty}{0.025 [0.007, 0.087]}  % probe_greedy_050
\newcommand{\SampledSeventyFive}{0.037 [0.013, 0.105]}  % probe_sampled_075
\newcommand{\GreedyGate}{0.362 [0.266, 0.472]}  % probe_greedy_gate
\newcommand{\SampledGate}{0.275 [0.189, 0.381]}  % probe_sampled_gate
\newcommand{\TsemAUC}{0.557 [0.309, 0.820]}  % trigger_tsem
\newcommand{\EpaAUC}{0.523 [0.321, 0.734]}  % trigger_epa
\newcommand{\PtwoCI}{-155 ms [-487, +59]}  % pass2
\newcommand{\SpecTTFAr}{+32 ms [-92, +92]}  % spec_ttfa
\newcommand{\SpecEnergyr}{+35.3 J [+34.6, +35.7]}  % spec_energy
\newcommand{\EnergyListenerr}{+17.8 J [+17.1, +18.6]}  % energy_listener
\newcommand{\AllocDecisions}{576}  % pass 6 decision records
\newcommand{\AllocBelow}{537}  % pass 6 decision records
\newcommand{\AllocInfeasible}{39}  % pass 6 decision records
\newcommand{\AllocSpends}{0}  % pass 6 decision records
\newcommand{\BreakEven}{0.1412}  % pass 6 decision records

\graphicspath{{figures/}}

\title{The Window Belongs to the Listener: Allocating Think-While-Listening
Compute on a Device That Pays for It}
\author{Ali Salimi Sadr}
\date{}

\begin{document}
\maketitle

% Abstract: written last, from RESULTS.md.

\section{Introduction}
\label{sec:intro}

People take turns with gaps of a few hundred milliseconds~\cite{stivers2009universals}.
A cascaded voice agent (recognizer, language model, synthesizer) running on an
edge device answers much later. On our device, a reactive pipeline answering
spoken multi-step math problems starts speaking a median of \TTFAMulti\,ms after
the user stops (Fig.~\ref{fig:stages}).

The field's answer is to \emph{think while listening}: start the language model
on partial transcripts, so that part of the reply is ready by the endpoint
\cite{predgen2025,ltsvoiceagent2026,epa2026,livemind2024,shanks2025,stitch2026,shih2025thinklisten}.
That answer assumes the reasoner is idle while the user speaks, and that the
recognizer runs somewhere else. On a unified-memory system-on-chip neither
holds. The recognizer, the reasoner and the synthesizer share one memory system
and one set of cores, and every speculative decode is paid for by whatever runs
next to it.

This paper measures what the silence is made of on such a device, and asks
which consumer the listening window should be spent on. The silence belongs to
the recognizer. Its final decode is the largest stage on every length we test:
\SttShareMulti\% of the stage sum on multi-step problems, against
\LlmShareMulti\% for the language model's first token (Fig.~\ref{fig:stages}).
Speculating the reply attacks the smallest stage. Speculating the
\emph{transcript} attacks the largest. If the recognizer commits the words two
consecutive hypotheses agree on while the user is still speaking, the final
decode at the endpoint shrinks to the uncommitted tail.

Both kinds of speculation cost something on shared silicon, so a rule has to
choose between them. Our contributions:
\begin{enumerate}
\item \textbf{The costs, measured on one axis.}
  \begin{itemize}
  \item A co-resident process taxes the recognizer by its mere residency,
    through transparent-huge-page fallback. A process that only holds pages
    adds between \TaxBallastA{} and \TaxBallastB{} to each final decode, in two
    runs.
  \item A draft still decoding at the endpoint delays the first audio by
    \ThinkerOverlap{} ($n=\ThinkerOverlapN$).
  \item A hypothesis still decoding at the endpoint costs the final
    \ListenerC\,ms per second of overlap.
  \end{itemize}
\item \textbf{A per-turn allocation rule.} A hypothesis or a draft is issued
  only if it can finish, and pay off, before the endpoint. The rule uses the
  measured usability of drafts and the expected commit time of hypotheses. It
  preempts an orphaned hypothesis at the endpoint by freezing it. On our task
  the rule never spends the window on the reasoner, and it spends it on the
  listener only where the utterance is long enough. Against a baseline that
  never speculates, it cuts the median paired time to first audio by
  \PthreeRed{} on multi-step problems ($n=\PthreeN$ items, pre-registered; the
  CI's lower end lies below the 300\,ms bar we registered). On short utterances
  it abstains, with no detectable cost (\DevTTFA).
\item \textbf{The record.} Every prediction was registered before its data
  existed, and every one is reported with its verdict, the falsified ones
  included (Table~\ref{tab:ledger}).
\end{enumerate}

\begin{figure}[t]
  \centering
  \includegraphics[width=\columnwidth]{paper_fig1_stage_decomposition.pdf}
  \caption{What the silence is made of: median stage durations after the end
  of speech, reactive pipeline without partial decodes, five utterance
  lengths.}
  \label{fig:stages}
\end{figure}

\section{Related work}
\label{sec:related}

\textbf{Speculating the reply.} Several systems start the language model
before the user has finished. They predict the rest of the input, trigger
generation on semantic or acoustic cues, or let the model reason over the
incoming stream
\cite{predgen2025,ltsvoiceagent2026,epa2026,livemind2024,ddtsr2026}. Spoken
language models fold listening and thinking into one model
\cite{shanks2025,stitch2026,miniomnireasoner2025,chronothink2026,shih2025thinklisten},
and full-duplex models remove the turn boundary altogether
\cite{moshi2024,lws2026}. Evaluations of turn-taking and voice agents
\cite{fullduplexbench2025,fullduplexbenchv3_2026,voicebench2026,voiceagentbench2025,evabench2026,provoicebench2026,sid2026}
and low-latency cascades \cite{ilava2025,chipchat2025,kame2026,streamrag2025,moshirag2026}
measure latency end to end. None of them prices speculation against the
recognizer on hardware the two share.

\textbf{Speculating the transcript.} Streaming recognition with local
agreement commits the prefix that consecutive hypotheses agree on
\cite{machacek2023whisperstreaming}, on top of Whisper
\cite{radford2023whisper}; others decide when a streaming recognizer should
emit \cite{x2streaming2026}. We use the mechanism, not invent it. Our
contribution is to put it head to head with reply speculation, under the same
issue rule and the same overlap penalty on shared silicon, with a controller
that chooses between them.

\textbf{Interference on shared SoCs.} Co-running neural workloads interfere
through shared memory bandwidth, caches and accelerators
\cite{haxconn2024,aliyun2018bwlock,kim2024interference,kang2026emcgovernor,sgdrc2025,quilt2026,prophet2017,igniter2022}.
We add a mechanism these models do not cover. A co-resident process taxes a
CPU decoder through huge-page fallback, without consuming bandwidth and without
running at all.

\section{Problem}
\label{sec:problem}

A turn ends when the voice-activity detector declares the endpoint, a fixed
settle time after the last speech. Let $D$ be the audio the recognizer holds at
the endpoint. Time to first audio (TTFA), from the true end of speech, is
\begin{equation}
\mathrm{TTFA} = s + F(D - C) + L + S,
\end{equation}
where $s$ is the settle time, $C$ the audio committed before the endpoint, $F$
the final decode of the uncommitted tail, $L$ the language model's
time to first token and $S$ the synthesizer's first chunk. The final is close
to affine in its audio, $F(a) = \phi_0 + \phi_1 a$. Two consumers can spend
the listening window.

\textbf{The listener} decodes the uncommitted buffer with a small model and
commits what two consecutive hypotheses agree on. It reduces $D - C$. A
hypothesis decoding when the endpoint arrives competes with the final for the
same cores. It adds $c$ milliseconds of final decode per second of overlap,
unless it is stopped.

\textbf{The thinker} drafts the reply from the transcript so far, at fraction
$f$ of the utterance. A draft is usable with probability $p_{\mathrm{usable}}(f)$,
and is worth its cost only above a break-even $p^\ast$. A draft still decoding
at the endpoint delays the reply by a fixed overlap penalty. A draft that
cannot finish before the endpoint is all cost. So the window it can use,
\begin{equation}
w(f) = D(1-f) - d(fD),
\end{equation}
the speech left after the decode $d$ of what has been heard, must exceed the
draft's own decode.

A rule therefore issues work only if it is \emph{feasible}: it can finish, and
for the listener commit, before the endpoint. It must also be \emph{valuable}:
for the thinker, $p_{\mathrm{usable}}(\hat f) \ge p^\ast$ at the estimated
fraction $\hat f$. Neither $D$ nor the endpoint is known in advance. The rule
works from a prior over utterance length and from live, in-pipeline timings.

\section{System}
\label{sec:system}

\textbf{Pipeline.}
\begin{itemize}
\item Silero voice-activity detection inside a Pipecat pipeline.
\item Two Whisper engines, through faster-whisper with int8 weights, on three
  shared CPU cores: \texttt{tiny} for hypotheses and \texttt{base} for the final.
\item A 4-bit Gemma~4~E2B served by llama.cpp with every layer on the GPU, on
  the other three cores, with one slot.
\item Piper for synthesis.
\end{itemize}
Every turn writes a record with its stage timestamps, its decodes and the power
rails.

\textbf{Committing while listening.} A hypothesis worker, a separate process,
decodes only the audio after the committed point. A LocalAgreement-2 committer
\cite{machacek2023whisperstreaming} commits the agreed prefix at segment
granularity, guarding the last 0.3\,s of the buffer. At the endpoint the final
decodes only the uncommitted tail.

\textbf{The issue rule.} Hypotheses are self-paced to a duty cycle. After a
decode of length $d$ the worker rests $(1/\delta - 1)\,d$, so the recognizer's
cores are never saturated and voice-activity detection is never starved.

\textbf{The feasibility gate.} A hypothesis is issued only if the decodes still
needed to commit its words, with their rests, fit in the expected remaining
speech. The expected remaining speech is a prior utterance length, derived at
run start as the median duration of the set's files, minus the speech heard so
far. On short utterances the gate abstains outright.

\textbf{Preemption at the endpoint.} A hypothesis still decoding at the
endpoint is frozen (\texttt{SIGSTOP}). It is resumed only once the recognizer
is idle after the reply has played, and its stale result is discarded by
sequence number. Freezing removes the listener's overlap penalty without
killing the worker. Two earlier designs did not survive their own validity
checks: killing and respawning it, and keeping a hot spare
(Table~\ref{tab:ledger}).

\textbf{The allocator.} For the thinker, at every partial transcript the
allocator estimates $\hat f$ from the prior. It looks up $p_{\mathrm{usable}}(\hat f)$
on a curve measured offline, and issues a draft of budget $B$ only if
$p_{\mathrm{usable}}(\hat f) \ge p^\ast$ and $B$ tokens fit in $w(\hat f)$ at the
live decode rate. Every decision is logged with its inputs and its reason
(Table~\ref{tab:allocator}).

\section{Experimental setup}
\label{sec:setup}

\textbf{Device.} NVIDIA Jetson Orin Nano Super, 8\,GB of unified memory, in
power mode MAXN\_SUPER with clocks pinned, run headless.

\textbf{Sets.} One benchmark, five utterance lengths: a 16-utterance
development set, two digit sets and the multi-step split of Spoken-MQA
\cite{spokenmqa}, whose problems derive from GSM8K \cite{cobbe2021gsm8k}. We
use 80 items drawn with a fixed seed, and as a fifth length the 20 longest of
them. Lengths are in Fig.~\ref{fig:stages}.

\textbf{Protocol.}
\begin{itemize}
\item Audio enters through the microphone's delivery path at real time, one
  file per turn. The next file starts only after the reply has finished and a
  fixed quiet period has passed, so turns never overlap.
\item TTFA is measured from the true end of speech, known to the file source,
  to the first reply audio.
\item Arms are compared on the same items. Each item's value is its median
  over three runs, the difference is taken per item, and the 95\% confidence
  interval is a percentile bootstrap over items.
\item Turns split by the voice-activity detector, turns whose endpoint fired
  after the final, and runs that fail a pre-registered validity check are
  excluded by rules fixed before scoring.
\end{itemize}

\textbf{Metrics.} TTFA and its stages, the word-error-rate difference against
the baseline on the same audio, energy per turn from the board's power rails
net of the idle floor, and the fraction of words committed at the endpoint.

\textbf{Pre-registration.} Every prediction, its bar and its scoring rule was
committed to the lab notebook before the run that tests it. A correction is
appended, never edited in. Table~\ref{tab:ledger} lists every prediction with
its registration commit and its verdict.

% Sections 6-8 (results, limitations, conclusion with the required
% generative-AI disclosure) follow from RESULTS.md.

\bibliographystyle{IEEEtran}
\bibliography{refs}

\end{document}
